% Fuente: recuento de funciones test_* en tests/ del repositorio.
\begin{tabularx}{\textwidth}{lrX}
\toprule
Fichero de contratos & Tests & Qué invariante protege \\
\midrule
\texttt{test\_economic\_contract.py} & 28 &
Que una venta respete costes, histéresis, mínimo de tenencia, suelo de diversificación y tope de
efectivo; que ninguna orden se emita sin destino mejor \\
\addlinespace
\texttt{test\_backtest\_contract.py} & 12 &
Que la contabilidad del \textit{backtest} no cobre costes de operaciones no ejecutadas ni resucite
posiciones vendidas \\
\addlinespace
\texttt{test\_catalog\_contract.py} & 10 &
Que el catálogo sea cerrado: toda variable declarada existe, toda rejilla es alcanzable y no hay
parámetros fuera de él \\
\addlinespace
\texttt{test\_statistical\_contract.py} & 8 &
Que las puertas de elegibilidad y el emparejamiento por periodo se comporten en sus límites, y que
un bloque incompleto se marque no aplicable \\
\addlinespace
\texttt{test\_workflow\_contract.py} & 7 &
Que los estados de ejecución, la persistencia de artefactos y la reconciliación de procesos muertos
sean consistentes \\
\addlinespace
\texttt{test\_alpha\_curve\_contract.py} & 6 &
Que la curva percentil--alfa produzca un valor distinto por rango, que una relación decreciente
caiga a la salvaguarda y que la ausencia de cohortes cerradas devuelva un valor no disponible \\
\addlinespace
\texttt{test\_identity\_contract.py} & 4 &
Que la identidad del \textit{dataset} y la clave de evaluación incluyan todo lo que las distingue,
incluido el hash de datos \\
\midrule
\textbf{Total} & \textbf{75} & \\
\bottomrule
\end{tabularx}
